\documentclass[../../../../SoftwareRequirementsSpecification.tex]{subfiles}
\begin{document}
% Richiede le macro definite in src/macros/useCaseMacros.tex
% (incluse nel preambolo di SoftwareRequirementsSpecification.tex)

\ucstart
\begin{xltabular}{\textwidth}{|l|X|}
  \hline
  \ucid{A-04} \\ \hline
  \ucname{Modifica Dati Anagrafici} \\ \hline
  \ucactor{Atleta} \\ \hline
  \ucdescription{L'atleta modifica i propri dati anagrafici
    (altezza, peso, data di nascita).} \\ \hline
  \ucprecond{L'atleta deve essere autenticato.} \\ \hline
  \ucpostcond{I dati anagrafici dell'atleta vengono
    aggiornati nel sistema.} \\ \hline
  \ucmainflow{
    \item L'atleta apre la pagina del proprio profilo.
    \item Il sistema mostra i dati anagrafici correnti.
    \item L'atleta preme il pulsante "Modifica Dati".
    \item Il sistema mostra il form di modifica dati (data di nascita,
      altezza, peso).
    \item L'atleta modifica i dati desiderati.
    \item L'atleta preme il pulsante "Salva Modifiche".
    \item Il sistema verifica e salva i dati nel database.
    \item Il sistema mostra un messaggio di conferma.
  } \\ \hline
  \ucaltflow{
    \textbf{7a. Dati non validi} \newline
    - Il sistema mostra un messaggio di errore ("Dati inseriti non
    validi"). \newline
    - Il flusso riprende dal passo 4.
  } \\ \hline
\end{xltabular}
\end{document}
